\documentclass[11pt,a4paper]{article}
\usepackage[margin=26mm]{geometry}
\usepackage{amsmath,amssymb,amsthm,booktabs,graphicx}
\usepackage{xcolor}
\usepackage[colorlinks=true,linkcolor=blue!45!black,citecolor=blue!45!black,urlcolor=blue!45!black]{hyperref}
\usepackage{microtype}
\setlength{\emergencystretch}{2em}
\newtheorem{proposition}{Proposition}
\title{From driven flow to inherited growth\\[.3em]\large Constructions and establishment tests}
\author{David Galbraith}
\date{Working draft, 11 September 2026}
\begin{document}
\maketitle
\begin{abstract}
Can variation in a driven physical system become inherited differential growth?
A complete answer requires a physical copying mechanism as well as a criterion
for its spread. We separate channel formation, flow sorting, persistent patterns,
reproduction with inherited differences, and portable information. Existing
lattice and reaction--diffusion constructions establish some of these capabilities
under specified rules; they do not yet realize the full sequence in one substrate.
We then derive a conditional establishment test from measured usable power,
copy cost, inheritance fidelity and the lifetime of copying ability. For a parent
that accumulates work before each attempt, the critical cost is
$E_c=W\log(1+p)/\mu$. Replacing that clock by Poisson attempts with the same
immortal-parent mean rate gives $E_c=Wp/\mu$, and can reverse the predicted outcome.
Independent life-history simulations check the difference. Both results apply
established branching and renewal mathematics. Their role is to make the proposed
physical connection testable. A common microscopic construction of variation,
copying and inheritance remains to be established.
\end{abstract}

\section{The claim to establish}

The long-term aim is to explain how local physical dynamics in an open, driven
system can produce inherited differences that affect reproduction. The present
claim is narrower: the necessary connections can be stated separately, some have
constructive examples, and copying-clock assumptions have measurable consequences
for establishment. A physical derivation must eventually join the connections
under the same interaction rules. An analogy between replication and an epidemic
does not establish that derivation.

The system exchanges matter or energy with maintained reservoirs. Let $J_i$ be
an intake throughput and $W_i$ its usable power after conversion losses. For a
specified driving potential, write $W_i=\epsilon_i J_i$, where $\epsilon_i$ is work
per unit throughput, not necessarily a dimensionless efficiency. A work budget
must also satisfy $\sum_i W_i\le P_{\mathrm{available}}$. Sustained currents require
external supply; the argument assumes no maximum-entropy-production principle.
The earlier concentration model has reversible coarse state dynamics despite
representing a maintained throughput. It is therefore not by itself a resolved
thermodynamic model of the driven apparatus.

\section{Five separate physical questions}

\begin{center}
\begin{tabular}{@{}p{.24\linewidth}p{.71\linewidth}@{}}
\toprule
Question & What would establish it?\\
\midrule
Paths form & Geometry develops under local material and transport rules.\\[3pt]
Flow sorts & A measured small fraction of available paths carries a specified
fraction of total throughput.\\[3pt]
A pattern persists & Spatial organisation continues while its constituent
material turns over.\\[3pt]
A pattern reproduces & An identifiable parent produces a new structure, with
perturbations that can persist across parent--offspring relations.\\[3pt]
Information travels & A transported state recreates an inherited property
at a separate site under common destination conditions.\\
\bottomrule
\end{tabular}
\end{center}

These are not necessarily a chronological ladder. Unequal shares among fixed
pipes establish sorting without circulation. A whirlpool can maintain its shape
as water passes through it, yet need not preserve a particular perturbation or
produce descendants. Spatially correlated patterns need not be related by copying.
Inheritance of copying ability also has a lifetime distinct from the lifetime
of the visible object. These distinctions determine which measurements matter.

The companion paper derives an exact stationary sorting window for a specified
conserved-flow diffusion \cite{sorting}. Independent contacts proportional to
actual flow recover its variance-proportional-to-flow case. Deterministic lattice
examples also seed unequal sharing from identical coarse states, without random
draws after initialization. Deterministic effective noise has mathematical
precedents \cite{melbourne}; the existing lattice's correlations and limiting
clock have not been fully matched to the theorem. This supplies a candidate
source of variation, not yet inherited variation.

\section{What the constructions establish}

The restored lattice-Boltzmann pole model produces returning trajectories in
its nearly steady wake, absent in low-drive, linear and no-pole controls.
The hydrodynamic connection of suitable lattice-Boltzmann rules to Navier--Stokes
is established \cite{lbm}. A material extension replaces the fixed pole mask by
binary sites with local cohesive exchanges. Their block-average occupancy sets
fluid reflection. Three prepared patches retain about 96\% spatial occupancy
overlap and have 30--43 returning tracers out of 120 in the frozen final field.
Equal material dispersed through the domain and uncoupled-material controls have
no returning tracers in that diagnostic. Removing cohesion spreads the patch
but leaves some circulation. The substrate supplies support; the experiment does
not derive a macroscopic patch from uniform noise, free suspension or erosion.

The repository also reproduces known Gray--Scott spot multiplication from a
finite seed using local reaction and diffusion laws \cite{pearson,spots}. It
contains no instruction to detect and copy a spot. The reaction network itself
supplies autocatalysis; related chemical spot replication has been observed
experimentally \cite{chemicalspots}. Splitting therefore has a constructive
precedent, but these runs do not establish selectable inherited differences.
Molecular templating provides another precedent for portable information simpler
than a complete genetic system \cite{peptide}; it has not emerged from the lattice
models here. Existing token-carrier simulations instead assume copying and expression.

A prepared rare configuration can test a possible mechanism without waiting for
its historical occurrence. It cannot measure its origin rate. To show that disturbance can create the capability, specify an initial ensemble
and an allowed perturbation history that reaches it under the same rules.
A prepared capable state can be the endpoint of that construction, but its
preparation alone does not establish reachability from the proposed background. The origin,
establishment and spread of a capability remain separate questions.

\section{From intake to establishment}

Consider a rare lineage against a fixed background. A copying-capable structure
receives constant usable power $W>0$ and starts with an empty store. Every attempt
costs $E>0$, resets that store to zero and produces one viable daughter carrying
the same copying capability with independent probability $p\in(0,1]$. The parent
survives a copying attempt. The daughter starts with an empty store. Copying ability
is lost irreversibly at exponential rate $\mu>0$; this may include physical death
and loss of the inherited state. Copies are independent once born. Failed attempts
still consume their work. Constant intake and independent descendants are rare-lineage
approximations, not a model of population takeover under a fixed total budget.

\begin{proposition}
For this accumulator model, positive probability of eventual lineage survival
occurs precisely when
\begin{equation}
 \mathcal R=\frac{p}{\exp(\mu E/W)-1}>1,
 \qquad E<E_c=\frac{W}{\mu}\log(1+p).
 \label{eq:clock}
\end{equation}
\end{proposition}
\begin{proof}
Attempts occur at ages $nt_c$, $n=1,2,\ldots$, where $t_c=E/W$. Survival of copying
ability to age $nt_c$ has probability $\exp(-\mu nt_c)$. The expected number of
viable daughters over a parent's life is therefore
$p\sum_{n\ge1}\exp(-\mu nt_c)$, giving Eq.~\eqref{eq:clock}. Identical independent
reproduction histories define a branching genealogy with nondegenerate offspring
law. Its survival threshold is mean offspring greater than one \cite{branching}.
Solving that inequality gives $E_c$. Birth ages affect growth in physical time;
they do not change this genealogical criterion.
\end{proof}

This is an application of established renewal and branching mathematics, not a
new general threshold theorem. More generally the mean offspring is the integral
of survival-weighted viable reproduction over age. Age-dependent reproduction
cannot in general be replaced by a single mean attempt rate \cite{age}.
For immediate offspring in this particular accumulator model, the renewal equation
also gives the mean-growth exponent $\lambda=W\log(1+p)/E-\mu$. No claim about
empirical growth or eventual population size follows without testing the model.

A parent with Poisson attempts at rate $W/E$ has the same mean attempt rate if
immortal, but a different lifetime output:
\begin{equation}
 \mathcal R_{\mathrm P}=\frac{Wp}{\mu E},\qquad
 E_{c,\mathrm P}=\frac{Wp}{\mu}.
 \label{eq:poisson}
\end{equation}
Because $\log(1+p)<p$, the Poisson replacement can predict establishment where
the accumulator predicts extinction. It permits attempts before enough work
has arrived in an empty store, so it is not an equivalent pathwise resource
model. At faithful copying ($p=1$), the two critical costs differ by a factor
$\log2$. This comparison identifies a model substitution that needs justification.

For example, take $W=\mu=1$, $p=.8$ and $E=.7$. The accumulator has
$\mathcal R=.789<1$; the Poisson model has $\mathcal R_{\mathrm P}=1.143>1$.
Thus a plausible change in the copying clock reverses the predicted outcome
without changing the nominal attempt rate. A visible structure can also be
long-lived while its copying state is short-lived: if these are independent
irreversible losses with rates $d$ and $r$, then $\mu=d+r$, not $d$ alone.

Portable copying can enter through the same quantities. If transit independently
destroys the required state at rate $\kappa$ over duration $\tau$, the surviving
fidelity becomes $p=p_0\exp(-\kappa\tau)$. Substitution in Eq.~\eqref{eq:clock}
gives a conditional transit-dependent critical cost. Transit also delays births
and changes the growth exponent; the immediate-offspring expression for $\lambda$
then no longer applies. The transport law and the copying mechanism must be
measured or constructed rather than inferred from a successful threshold fit.

For multiple inherited states with Markov transitions $T$, disappearance rates
$d$, exponential attempt rates $b$, and viable daughter-state matrix $P$, the
corresponding known criterion uses
$K=(\operatorname{diag}d-T)^{-1}\operatorname{diag}(b)P$ and spectral radius
$\rho(K)>1$, under finite irreducible branching assumptions \cite{branching}.
The repository checks this conditional model independently. Its rate-only state
space cannot represent an empty-store accumulator without adding age or stored
work. The scalar example shows why that distinction can change the conclusion.

\section{Numerical checks and physical tests}

Independent event histories check the lifetime output in Eqs.~\eqref{eq:clock}
and \eqref{eq:poisson}. Each setting uses 100,000 parents; a sampled loss time
terminates explicit regular or exponential attempt times. Bernoulli outcomes
record viable daughters. These are individual reproduction histories, not physical
observations or complete population simulations.

\begin{center}
\begin{tabular}{@{}lrrr@{}}
\toprule
Clock & Cost & Predicted daughters & Observed mean (SE)\\
\midrule
Accumulator & .45 & 1.4077 & 1.4006 (.0058)\\
Accumulator & .70 & .7891 & .7836 (.0037)\\
Accumulator & 1.00 & .4656 & .4643 (.0026)\\
Poisson & .45 & 1.7778 & 1.7643 (.0070)\\
Poisson & .70 & 1.1429 & 1.1355 (.0049)\\
Poisson & 1.00 & .8000 & .7960 (.0038)\\
\bottomrule
\end{tabular}
\end{center}

All six means lie within two estimated standard errors of their predictions;
this is a descriptive numerical check, not a simultaneous confidence statement.
The accumulator never spends work before receipt in the retained runs. Three
unit tests check the survival-weighted sum, the classification reversal and
work accounting. Raw outcomes, seeds and source hashes remain in the repository.
Existing spatial, carrier and multitype checks are retained separately.

A discriminating experiment would measure usable intake, loss of copying ability,
copy cost and parent--daughter fidelity independently of lineage outcomes.
Time-resolved attempts distinguish accumulation from an exponential clock.
Varying intake or cost then tests the predicted crossing of one viable daughter
per lifetime; replicate genealogies assess establishment separately. Individual
history and perturbation tracking must distinguish inheritance from common
environmental forcing. More total copies are not sufficient if they systematically
lose the copying state. Fixed-background predictions apply only while the lineage
is rare; continued growth requires explicit accounting for competition and losses.

The strongest next construction would let a flow-maintained pattern acquire a
copying capability through a finite perturbation, then transmit a modified
property to descendants under unchanged local interaction rules. Parent and
recipient manipulations should separate copied state from shared surroundings.
A transported carrier would need to recreate that property at a new site.
Parameters measured in those histories could then predict lineage outcomes on
held-out runs. This would join physical possibility to a quantitative test;
programming the copying operation would leave the central connection assumed.

\section{Scope of the proposed contribution}

The ambition is a physical account of inherited differential growth in driven
systems. The constructions support individual capabilities; the copying-clock
comparison states one conditional prediction in independently measurable terms.
Together they do not yet establish a common microscopic origin of variation,
copying and portable inheritance. The decisive advance would be a construction
that joins those capabilities under the same rules and predicts lineage outcomes
from separately measured histories, or a physical test that distinguishes the
proposed mechanism from alternatives. The present result gives that programme
a concrete test while keeping its assumptions separate from its ambition.

\paragraph{Reproducibility.}
Sources, retained outcomes, assumptions and test commands are at
\url{https://github.com/divadwg/selection-as-physics}, with the two-paper map in
\texttt{paper/README.md}. The companion sorting paper is independent of this
paper's copying assumptions.

\begin{thebibliography}{99}
\bibitem{sorting} D. Galbraith, A noise-scaling window for concentration of a shared flow,
companion working paper (2026), repository \texttt{paper/finite\_noise\_window.pdf}.
\bibitem{melbourne} I. Melbourne and A. M. Stuart, A note on diffusion limits of chaotic
skew-product flows, \emph{Nonlinearity} \textbf{24}, 1361--1367 (2011), corrected version (2015).
\href{https://arxiv.org/abs/1101.3087}{arXiv:1101.3087}.
\bibitem{lbm} M. Junk, A. Klar and L.-S. Luo, Asymptotic analysis of the lattice Boltzmann equation,
\emph{Journal of Computational Physics} \textbf{210}, 676--704 (2005).
\href{https://doi.org/10.1016/j.jcp.2005.05.003}{doi:10.1016/j.jcp.2005.05.003}.
\bibitem{pearson} J. E. Pearson, Complex patterns in a simple system,
\emph{Science} \textbf{261}, 189--192 (1993).
\href{https://arxiv.org/abs/patt-sol/9304003}{arXiv:patt-sol/9304003}.
\bibitem{spots} W. N. Reynolds, S. Ponce-Dawson and J. E. Pearson,
Self-replicating spots in reaction-diffusion systems,
\emph{Physical Review E} \textbf{56}, 185--198 (1997).
\href{https://doi.org/10.1103/PhysRevE.56.185}{doi:10.1103/PhysRevE.56.185}.
\bibitem{chemicalspots} K.-J. Lee, W. D. McCormick, J. E. Pearson and H. L. Swinney,
Experimental observation of self-replicating spots in a reaction--diffusion system,
\emph{Nature} \textbf{369}, 215--218 (1994).
\href{https://doi.org/10.1038/369215a0}{doi:10.1038/369215a0}.
\bibitem{peptide} D. H. Lee et al., A self-replicating peptide,
\emph{Nature} \textbf{382}, 525--528 (1996).
\href{https://doi.org/10.1038/382525a0}{doi:10.1038/382525a0}.
\bibitem{branching} S. Hautphenne, G. Latouche and G. Nguyen,
Extinction probabilities of branching processes with countably infinitely many types,
\emph{Advances in Applied Probability} \textbf{45}, 1068--1082 (2013).
\href{https://arxiv.org/abs/1211.4129}{arXiv:1211.4129}.
\bibitem{age} E. B. Stukalin, I. Aifuwa, J. S. Kim, D. Wirtz and S. X. Sun,
Age-dependent stochastic models for understanding population fluctuations in
continuously cultured cells, \emph{J. R. Soc. Interface} \textbf{10}, 20130325 (2013).
\href{https://doi.org/10.1098/rsif.2013.0325}{doi:10.1098/rsif.2013.0325}.
\end{thebibliography}
\end{document}
